\documentclass[11pt]{article}
\usepackage{amsmath,amssymb,amsfonts}
\usepackage{geometry}
\usepackage{enumitem}
\usepackage{color}
\usepackage{hyperref}
\usepackage{booktabs}

\geometry{letterpaper, margin=1in}

\definecolor{addressed}{rgb}{0,0.5,0}
\definecolor{partial}{rgb}{0.8,0.4,0}
\definecolor{pending}{rgb}{0.8,0,0}

\title{\textbf{Response to Reviewers --- Round 2}\\[0.5em]
\large Manuscript: Flux-Preserving Adaptive Finite State Projection\\
for Multiscale Stochastic Reaction Networks\\[0.5em]
\normalsize Submitted to: SIAM Multiscale Modeling and Simulation}

\author{Aditya Dendukuri, Shivkumar Chandrasekaran, and Linda Petzold\\
Department of Computer Science\\
University of California, Santa Barbara}

\date{\today}

\begin{document}

\maketitle

\noindent Dear Editor and Reviewers,

\medskip

Thank you for the continued careful evaluation of our manuscript and for the
constructive feedback in this second round. We are grateful for the referees'
thorough reading, and we have addressed every point raised. Below we respond
point by point.

%======================================================================
\section*{Response to Referee 1}
%======================================================================

\noindent\textbf{Comment 1.1:} \textit{The ``total flux'' of a reduced state
space is defined as the sum of the flux of each state in the subset. The
majority of this total flux would then be the propensities to move between two
states already contained in the reduced state space. If the goal of the
algorithm is to avoid loss of states that move to other metastable regions,
would it make more sense to limit the flux to propensities of moving outside
the reduced state space?}

\textbf{Response:} This is an excellent question that cuts to the heart of the
design choice.  One could protect states based on their \emph{boundary outflux}
$\Phi_{\partial}(\boldsymbol{x},t):=p(\boldsymbol{x},t)\sum_{\boldsymbol{y}\in J'}\mathbf{A}_{\boldsymbol{y}\boldsymbol{x}}$ (the rate at which probability leaves $\boldsymbol{x}$ to states outside the current active set $J$), rather than total flux $\Phi(\boldsymbol{x},t)=p(\boldsymbol{x},t)\,w(\boldsymbol{x})$.
Total flux is preferable for two reasons, most clearly illustrated by the
bottleneck example.

\textbf{Boundary outflux fails at the mature stage of the computation.}
When $s_b$ first enters the active set, its destination states are not yet in $J$, so $\Phi_{\partial}(s_b,t)$ is large and boundary-outflux protection would correctly retain it.
However, once the destination regions are brought into $J$ by expansion (they are high-probability states that survive the quantile filter), $s_b$'s transitions lead entirely within $J$ and its boundary outflux drops to approximately zero.
At that mature stage, boundary-outflux protection would \emph{not} retain $s_b$, even though pruning it now disconnects the two regions.
Total flux $p(s_b)\,w(s_b)$ is large throughout \emph{both} stages because $w(s_b)$ is intrinsically large: it measures the rate at which probability cycles \emph{through} $s_b$ regardless of whether the destinations are in $J$ or $J'$.
A large exit rate combined with small occupancy is precisely the hallmark of a fast-transit connector, and total flux correctly identifies it at every stage of the computation.

\textbf{Circular dependency.}
Computing $\Phi_{\partial}(\boldsymbol{x},t)$ requires knowing $J'$, which is itself determined by the pruning decision.
Total flux $p(\boldsymbol{x})\cdot(-\mathbf{A}_{\boldsymbol{x}\boldsymbol{x}})$ is a byproduct of the generator diagonal and requires no knowledge of the pruning outcome.

We added Remark~3.1 to the revised manuscript making this distinction explicit.

\medskip

\noindent\textbf{Comment 1.2 (minor):} \textit{Line 140: the term for flux
$w(x)$ is used before being defined.}

\textbf{Response:} Thank you for catching this.  We added an inline definition
of $w(x)$ at its first use in Remark~2.5, so the reader does not need to look
ahead to Section~3 for the formal definition.

\medskip

\noindent\textbf{Comment 1.3 (minor):} \textit{Fig.\ 2: As the flux is central
to your paper, it would be helpful to visualize it for the readers. A natural
place would be in Fig.\ 2A, where the fluxes can be illustrated on the arrows.}

\textbf{Response:} We appreciate this suggestion and updated Figure~2A to
visualize probability flux on the arrows.
Transitions incident on the bottleneck state $s_b$ are drawn with thicker lines to indicate that they carry high flux $\Phi(s_b,t)=p(s_b,t)\,w(s_b)$, while the internal metastable transitions (within regions $A$ and $B$) are drawn thin to reflect their comparatively low flux.
The label $\Phi(s_b)=p(s_b)w(s_b)$ appears on the representative $s_1\leftrightarrow s_b$ arrow, and the annotation $p(s_b)\!\ll\!1$, $w(s_b)\!\gg\!0$ below $s_b$ makes explicit the key property: $s_b$ is rarely occupied yet transitions through it at a high rate, so its flux is non-negligible despite its low probability.
This directly illustrates why probability-only pruning fails while flux-based pruning retains $s_b$.

%======================================================================
\section*{Response to Referee 2}
%======================================================================

\subsection*{Theoretical Issues}

\noindent\textbf{Comment 2.1:} \textit{There is a significant notational error
in Eq.\ (4.4) regarding the definition of boundary fluxes. Given the stated
column convention, $\Phi_{\mathrm{out}}$ should involve $\mathbf{A}_{J'J}$
rather than $\mathbf{A}_{JJ'}$.}

\textbf{Response:} We thank the referee for catching this error.  We added an
explicit clarification of the column convention immediately after Eq.~(4.4):
the block $\mathbf{A}_{J'J}$ has rows indexed by $J'$ and columns indexed
by~$J$, so the product $\mathbf{A}_{J'J}\boldsymbol{p}_J$ gives the rate
at which probability flows \emph{out} of~$J$ into~$J'$. This convention is now
stated clearly in the text, and the notation is consistent throughout Section~4.

\medskip

\noindent\textbf{Comment 2.2:} \textit{The derivation for the adaptive
time-step rule ($\Delta t = \varepsilon_{\Delta t}/\Phi_{\max}$) relies on the
inequality $\Phi_{\mathrm{in}} + \Phi_{\mathrm{out}} \leq \Phi_{\max}$, for
which I fail to see the justification without a multiplicative constant that
accounts for the boundary's size or structure.}

\textbf{Response:} We redesigned the time-stepping
rule from first principles to remove the unjustified inequality.

The new rule $\Delta t_n = \varepsilon_{\Delta t}/\Phi_{\mathrm{total}}(J_n,t_n)$, where $\Phi_{\mathrm{total}} := \sum_{\boldsymbol{x}\in J_n}p(\boldsymbol{x},t_n)\,w(\boldsymbol{x})$, is justified in Proposition~3.5 by a semigroup expansion argument on the compressed dynamics. Expanding $e^{\widetilde{\mathbf{A}}_{JJ}\Delta t}\widetilde{\boldsymbol{p}}_n$ to first order gives
\[
\|\widehat{\boldsymbol{p}}(t+\Delta t)-\widetilde{\boldsymbol{p}}_n\|_1 \leq 2\,\Phi_{\mathrm{total}}\,\Delta t + O(\Delta t^2),
\]
so the total variation distance, and hence the leading-order fraction of probability mass redistributed within the active set, is bounded by $\Phi_{\mathrm{total}}\,\Delta t + O(\Delta t^2)$. The choice $\Delta t = \varepsilon_{\Delta t}/\Phi_{\mathrm{total}}$ therefore bounds this fraction by $\varepsilon_{\Delta t} + O(\Delta t^2)$, with an explicit constant.

Crucially, $\Phi_{\mathrm{total}} = \mathbb{E}_p[w(\boldsymbol{x})]$ is the probability-weighted mean propensity, which tracks the \emph{active} timescale of the system rather than the worst-case propensity. When probability resides in low-propensity states, $\Phi_{\mathrm{total}}$ is small and the step is large; when probability shifts into high-propensity states, $\Phi_{\mathrm{total}}$ grows and the step shrinks automatically. This self-adaptation, proved in Proposition~3.5, is what makes the rule effective across the stiff Robertson and oscillatory Oregonator benchmarks without manual tuning.

We moved the time-step proposition to Section~3 (Adaptive Time Stepping), separate from the error analysis, since it is a temporal resolution criterion rather than an error bound. Section~4 then simply invokes $\Phi_{\mathrm{total}}\,\Delta t_n = \varepsilon_{\Delta t}$ as a consequence of the step rule.

\medskip

\noindent\textbf{Comment 2.3:} \textit{The algorithm (3.1 and 3.3) performs
renormalization after pruning, yet the theoretical analysis in Section~4 is
built upon a compressed (reflecting) generator. The relationship between these
two treatments needs explicit clarification to ensure the error bounds remain
valid for the implemented code.}

\textbf{Response:} This is an important point, and we are happy to clarify:
there is \emph{no gap} between the theory and the implementation.

\textbf{The implementation constructs the compressed generator directly.}
Algorithm~3.2 (forward-enumeration matrix construction) builds the generator column by column. For each source state $x \in J$ and each reaction~$k$, the propensity $\alpha_k(x)$ contributes to the off-diagonal entry at row $y = x + \boldsymbol{\nu}_k$ \emph{only if} $y \in J$. The diagonal entry for column~$x$ is the negative of the sum of included off-diagonal entries. Transitions whose destination falls outside~$J$ are excluded from both the off-diagonal and the diagonal, so column sums are exactly zero by construction, which is precisely the defining property of the compressed generator $\widetilde{\mathbf{A}}_{JJ}$ (Proposition~4.1, property~1). The evolution step $e^{\widetilde{\mathbf{A}}_{JJ}\Delta t}\boldsymbol{p}$ therefore conserves total probability exactly.

\textbf{The renormalization addresses pruning losses, not generator leakage.}
The renormalization $\boldsymbol{p} \gets \boldsymbol{p}/\|\boldsymbol{p}\|_1$ in Algorithm~3.1 (Stage~3) compensates for the probability mass discarded when low-probability states are removed during pruning. It has nothing to do with the generator. Let $m \leq \alpha$ be the pruned mass. The $\ell^1$ distance between the original distribution and the renormalized one is
\[
\|\boldsymbol{p} - \widetilde{\boldsymbol{p}}_n\|_1 = m + \frac{m}{1-m}(1-m) = 2m \leq 2\alpha,
\]
where the first term is the mass removed by pruning and the second is the inflation of the retained entries by renormalization. This $2\alpha$ term appears in the local error bound (Corollary~4.6) and accumulates to $2N\alpha$ in the global bound (Corollary~4.7). The error analysis in Section~4, formulated entirely in terms of $\widetilde{\mathbf{A}}_{JJ}$, therefore applies directly to the implemented algorithm without any additional approximation gap.

\medskip

\noindent\textbf{Comment 2.4:} \textit{Estimating $w_{\max}$: Corollary~4.7
uses $w_{\max}$ (the maximum exit rate of pruned states), but it is unclear how
this value is robustly estimated on a truncated domain since it involves states
outside the current active set.}

\textbf{Response:} \textbf{Note on renumbering.}  The global error section has been substantially restructured in this revision.  The previous Corollary~4.7 (local error bound with pruning tolerances), which the referee is referring to, contained a proof gap: the claimed bound $\Phi_{\mathrm{out}} \leq \varepsilon_{\mathrm{flux}}\cdot\Phi_{\mathrm{total}}$ carried an implicit factor of $|J'|$ that was not justified.  We replaced it with two corollaries that have rigorous proofs: the new \textbf{Corollary~4.6} (Local step error) bounds the one-step error using operator norms and is where $w_{\max}(J'_n)$ now appears; the new \textbf{Corollary~4.7} (Global error bound) gives the accumulated error over $N$ steps and no longer contains $w_{\max}$ directly (the influx-from-pruned-mass effect is captured by the prior-error term via the CME contraction).  Note that $\varepsilon_{\mathrm{flux}}$ does not appear directly in either corollary; its role is qualitative — it controls which states end up in $J'_n$ and thereby keeps $w_{\max}(J'_n)$ small in practice. This is stated explicitly in the paragraph introducing Corollary~4.6 in Section~4.2.

We clarified the $w_{\max}$ estimation question in Corollary~4.6.  The exit rate $w(x) = \sum_k \alpha_k(x)$ is computed for every state $x$ in the active set \emph{before} pruning takes place (it is needed for the flux computation in Algorithm~3.1, Stage~2).  Since pruning only removes states already in the active set, $w_{\max}$ over the pruned states is available as a byproduct of the flux computation, with no need to evaluate $w(x)$ at states outside the active set.  Corollary~4.6 now states this explicitly.

\subsection*{Empirical Issues}

\noindent\textbf{Comment 2.5:} \textit{The manuscript lacks head-to-head
comparisons against strong adaptive FSP baselines, such as ETFSP or other
time-varying variants.}

\textbf{Response:} We added Section~5.6, which provides a head-to-head
comparison of FP-FSP against the Krylov-FSP-SSA algorithm of Vo and
Sidje~[Sidje2015]. This baseline shares
the same expand--evolve--prune skeleton and Krylov matrix exponential
integrator, differing in pruning (joint probability-and-derivative threshold per
Algorithm~2 of~[Dinh2016], without flux-based protection), time stepping
(mass-conservation accept/reject rather than flux-adaptive), and expansion
(SSA-driven with reachability enumeration). The comparison isolates the effect
of flux-aware pruning and flux-adaptive time stepping on three benchmark systems.

On the Oregonator, flux-adaptive time stepping provides a decisive advantage:
FP-FSP completes the simulation in 36\,s with $|S| \leq 1{,}057$, while
Krylov-FSP-SSA reaches only $t = 0.003$ (2.7\% of the time horizon) after
3{,}000 iterations and 74\,s, with the state space growing to 13{,}294 states.
On the bottleneck system, both methods correctly retain the connector state
$B{=}1$, but Krylov-FSP-SSA produces $\langle C \rangle = 0.12$ (40$\times$
below the reference $\langle C \rangle = 4.978$) because its step size grows
exponentially ($\tau : 100 \to 3{,}700$\,s) and the depth-$r{=}1$ expansion
cannot track the probability wavefront; FP-FSP avoids this because its
flux-adaptive step $\delta t = \varepsilon_{\Delta t}/\Phi_{\mathrm{total}}$
naturally limits to $\sim 1/k_2 = 10$\,s in the bottleneck phase, keeping
expansion pace-matched with the wavefront. On the Robertson system over $[0, 10^4]$,
only FP-FSP completes the horizon: the stiffness of $k_2 = 3\times10^7$ locks
Krylov-FSP-SSA's step size in $[3\times10^{-5},\, 3\times10^{-4}]$, reaching
only $t = 228$ (2.3\%) in 500 iterations.

These results are presented in Figure~7 of the revised manuscript.

\medskip

\noindent\textbf{Comment 2.6:} \textit{The authors should attempt to
numerically demonstrate the specific benefits of flux-based protection versus
pure probability pruning, and of $\Phi_{\max}$-based stepping versus other
activity indicators.}

\textbf{Response:} The new Section~5.6 directly demonstrates both benefits.

\textbf{Flux-based protection versus pure probability pruning.}
The bottleneck system provides the clearest illustration: with pure probability pruning (no flux criterion), the connector state $s_b$ ($B{=}1$) is removed because its occupancy $p(s_b)$ is small, severing the only path between the two metastable regions and causing the simulation to fail entirely.
FP-FSP retains $s_b$ because its flux $p(s_b)\,w(s_b)$ is large, yielding $\langle C \rangle = 4.978$ ($<0.1\%$ error).
On the Oregonator, flux-based pruning also keeps the active set compact ($|S| \leq 1{,}057$): the flux criterion identifies and discards dynamically irrelevant states, whereas the probability-and-derivative filter used by Krylov-FSP-SSA retains all states with significant positive derivative, causing the projection to grow to 13{,}294 states without bound.

\textbf{$\Phi_{\mathrm{total}}$-based stepping versus other activity indicators.}
$\Phi_{\mathrm{total}}$-based stepping adjusts $\Delta t$ from $1.6\times10^{-6}$ to $1.9\times10^{-5}$ on the Oregonator to match instantaneous system activity, while the accept/reject mechanism of Krylov-FSP-SSA locks at $\tau \approx10^{-6}$ because large propensities cause mass loss for any larger step.
On the bottleneck system, the flux-adaptive step naturally limits to $\sim 1/k_2$, keeping expansion pace-matched with the probability wavefront, while Krylov-FSP-SSA's exponentially growing $\tau$ produces a 14-state projection that misses the wavefront entirely, yielding $\langle C \rangle = 0.12$ (40$\times$ error).
Quantitative results are shown in Figure~7.

\medskip

\noindent\textbf{Comment 2.7:} \textit{Some guidance for choosing the key
tolerances $(\alpha, \varepsilon_{\mathrm{flux}}, \varepsilon_{\Delta t})$
across different types of models would be very valuable.}

\textbf{Response:} We agree this is important for practitioners and added a
parameter guidance paragraph at the start of Section~5.  Three practical rules
emerge from the benchmark experience.

\begin{itemize}
  \item \textbf{$\varepsilon_{\Delta t}$:} Well-mixed systems (Oregonator, Robertson, toggle switch) use $\varepsilon_{\Delta t} = 1$ because $\Phi_{\mathrm{total}}$ already scales with the dominant reaction rates and this keeps the per-step activity fraction bounded.
  Reduce $\varepsilon_{\Delta t}$ to $0.01$--$0.1$ when a long integration horizon demands higher accuracy, since local errors accumulate additively over $N$ steps (Corollary~4.7).

  \item \textbf{$\alpha$:} For connectivity-critical systems, $\alpha$ can be set aggressively ($0.7$--$0.9$) because the flux criterion rescues any misclassified connector states; the quantile filter then removes only the genuinely irrelevant tail.
  Without bottlenecks, moderate values ($0.1$--$0.3$) avoid over-pruning.

  \item \textbf{$\varepsilon_{\mathrm{flux}}$:} Set to $1$ for oscillatory or well-mixed systems without bottlenecks.
  Since no individual state can carry flux equal to the total system flux $\Phi_{\mathrm{total}}$, the protection threshold $\varepsilon_{\mathrm{flux}}\cdot\Phi_{\mathrm{total}}$ is unreachable, effectively disabling flux-based protection and letting the quantile filter operate freely.
  For bottleneck systems, match $\varepsilon_{\mathrm{flux}}$ to the expected relative flux of the connector state: $\varepsilon_{\mathrm{flux}}\sim p(s_b)\,w(s_b)/\Phi_{\mathrm{total}}$; in practice $10^{-12}$--$10^{-6}$ covers most cases.
\end{itemize}

%======================================================================
\section*{Summary}
%======================================================================

We sincerely thank both referees for their continued careful and constructive
reading. All outstanding concerns have been addressed.

On the theoretical side, we corrected the notational error in Eq.~(4.4) flagged
by Referee~2. We redesigned the time-step rule from first principles: Proposition~3.5
derives $\Delta t_n = \varepsilon_{\Delta t}/\Phi_{\mathrm{total}}$ from the CME
semigroup expansion, as detailed in the response to Comment~2.2. We clarified the
relationship between the implemented renormalization and the compressed generator
in the response to Comment~2.3. We substantially restructured the error analysis in Sections~4.2--4.3: the previous Corollary~4.7 contained a proof gap (the $\varepsilon_{\mathrm{flux}}\cdot\Phi_{\mathrm{total}}$ bound on $\Phi_{\mathrm{out}}$ was unjustified), which we corrected by replacing it with a rigorous operator-norm argument.  The result is split across the new Corollary~4.6 (Local step error, now in Section~4.2, which contains $w_{\max}$) and Corollary~4.7 (Global error bound, Section~4.3).  $w_{\max}$ estimation is addressed in Corollary~4.6, as explained in the response to Comment~2.4.

On the empirical side, the new Section~5.6 provides a head-to-head comparison
of FP-FSP against Krylov-FSP-SSA on three benchmark systems, directly
demonstrating the benefits of flux-based pruning (bottleneck system) and
flux-adaptive time stepping (Oregonator).

In response to Referee~1's conceptual question, Remark~3.1 explains why total
flux is the correct protection criterion: boundary outflux is approximately zero
for genuine bottleneck states (because their transitions lead into already-retained
high-probability regions) and is circular to compute during the pruning step.

Finally, the parameter guidance paragraph at the start of Section~5 addresses
Referee~2's request, giving concrete recommended ranges for
$(\alpha,\varepsilon_{\mathrm{flux}},\varepsilon_{\Delta t})$ across the four
benchmark model types, grounded in the theoretical error bounds and the benchmark
experiments.

Beyond the above, the manuscript includes several other revisions.  The
adaptive time-step rule (previously embedded in the error analysis) has been
moved to its own subsection (Section~3.4) to clarify that it is a resolution
criterion rather than an error bound.  Section~2.3 (\emph{Flux and Network
Connectivity}) has been expanded with a new figure (Fig.~3) showing how
removing the bottleneck state $s_b$ renders the generator exactly
block-diagonal, directly motivating the flux-protection criterion.  A Code
Availability section has been added pointing to the public repository
(\url{https://github.com/AdityaDendukuri/DiscStochSim.jl}), which contains
the full implementation and reproducible scripts for all figures, including
the Krylov-FSP-SSA comparison.

\vspace{1em}
\noindent Sincerely,\\[0.5em]
Aditya Dendukuri, Shivkumar Chandrasekaran, and Linda Petzold

\end{document}
